%==============================================================
%  Seonghwan Lim - Resume
%  Compile:  pdflatex SeonghwanLim_Resume.tex
%==============================================================
\documentclass[letterpaper,10pt]{extarticle}

\usepackage[left=0.45in,right=0.45in,top=0.34in,bottom=0.34in]{geometry}
\usepackage{cmap}
\usepackage[T1]{fontenc}
\usepackage[utf8]{inputenc}
\usepackage{charter}
\usepackage{titlesec}
\usepackage{enumitem}
\usepackage{xcolor}
\usepackage[hidelinks]{hyperref}

\definecolor{ink}{HTML}{1A1C20}
\definecolor{rule}{HTML}{1A1C20}

\pagestyle{empty}
\setlength{\parindent}{0pt}
\linespread{0.93}\selectfont
\color{ink}

% ---------- section headings ----------
\titleformat{\section}
  {\normalfont\bfseries\footnotesize}
  {}{0em}{\MakeUppercase}[\titlerule]
\titlespacing*{\section}{0pt}{4.5pt}{2pt}

% ---------- bullets ----------
\setlist[itemize]{
  leftmargin=11pt, topsep=1pt, itemsep=0.8pt, parsep=0pt,
  label=-
}

% ---------- role heading: role + employer (left) / dates (right), then project ----------
\newcommand{\role}[3]{%
  \textbf{#1}\hfill #3\par
  {\small\itshape #2}\par\vspace{0.5pt}
}

\begin{document}

%==============================================================
% HEADER
%==============================================================
\begin{center}
  {\LARGE\bfseries Seonghwan Lim}\\[2pt]
  {\small Computer Vision Engineer}\\[3pt]
  {\footnotesize
    Champaign, IL \textbar\ +1 (312) 287-3570 \textbar\
    \href{mailto:qube20101219@gmail.com}{qube20101219@gmail.com} \textbar\
    \href{https://seonghwan97.github.io}{seonghwan97.github.io}
  }
\end{center}

%==============================================================
% SUMMARY
%==============================================================
\section{Summary}
Computer vision engineer with an M.S. in Computer Science and experience developing, evaluating, and
optimizing end-to-end computer vision and machine learning systems. My work spans data preparation,
model development, performance analysis, and validation.

%==============================================================
% EDUCATION
%==============================================================
\section{Education}
\textbf{Illinois Institute of Technology} --- M.S. Computer Science, May 2026, GPA 3.9/4.0\\
\textbf{Seoul Nat'l University of Science and Technology} --- B.S. Electrical and Information
Engineering, Aug 2023, GPA 3.6/4.0

%==============================================================
% TECHNICAL SKILLS
%==============================================================
\section{Technical Skills}
\textbf{Languages:} Python, C/C++, SQL, Bash\\
\textbf{Computer Vision:} Object detection, classification, multi-object tracking, 2D/3D pose estimation, image processing, OpenCV\\
\textbf{Machine Learning:} PyTorch, CNNs, Vision Transformers, multi-task learning, data augmentation, scikit-learn, XGBoost\\
\textbf{Data and Tools:} Dataset design, labeling QC, NumPy, Pandas, Git, Linux, Docker, real-time GPU inference (PyTorch, CUDA), embedded systems

%==============================================================
% EXPERIENCE
%==============================================================
\section{Experience}

\role{Research Assistant, iConSenSe Lab, Illinois Institute of Technology}{Camera-only worker safety monitoring; mixed-reality safety training analysis}{Jan 2025 - Present}
\begin{itemize}
  \item Built a camera-only worker monitoring system on ordinary video, no wearables. A YOLO pose
        model extracts keypoints, ByteTrack holds worker identity across frames, and downstream heads
        roll activity, work state, and posture into a 10-second exposure score at a median 23 fps for
        3 to 9 workers.
  \item Checked every module against an independent reference across 42 subjects: 0.861 activity
        accuracy, 0.775 work-state F1, and 0.793 correlation with a 3D-pose reference. Scoped workload
        as an exposure indicator rather than a fatigue measurement after it reached only 0.526 against
        indirect calorimetry, cross-checked in the field against a CO2 sensing rig I built on
        Raspberry Pi and Jetson.
  \item On a separate mixed-reality study, derived movement and energy-expenditure features from
        single-camera video across 61 HoloLens 2 sessions using 2D pose estimation, MotionBERT 3D
        lifting, and spatiotemporal feature extraction, with no depth camera or motion capture rig.
\end{itemize}

\role{Research Assistant, MMCom Lab, Illinois Institute of Technology}{Fine-grained road surface classifier}{Jan 2025 - Jul 2026}
\begin{itemize}
  \item Trained a classifier over 27 road-surface types on RSCD, roughly 960k training images
        against a 50k test split. Traced most errors to a long-tailed class distribution and to
        surface types that look nearly identical in a single frame.
  \item Fixed the data before the model, filling out rare classes with diffusion-generated images,
        then tried multi-task and mixture-of-experts heads to separate the look-alike classes, moving
        EfficientNet-B0 from 91.70\% to 93.33\%.
\end{itemize}

\role{Research Assistant, Bhushan Research Group, Illinois Institute of Technology}{Image-based process control for microfluidic mixing}{Jan 2025 - May 2026}
\begin{itemize}
  \item Wrote the measurement code that registers and normalizes CFD contour and experimental
        fluorescence images, then computes a 0 to 1 mixing index from outlet pixel intensity.
  \item Benchmarked SVR, MLP, XGBoost, Random Forest, and linear regression as CFD surrogates. SVR
        won on held-out data (MAE 0.0014, R2 0.9998), close enough to replace a full CFD solve during
        inverse screening, confirmed with fresh CFD runs and six fluorescence experiments.
\end{itemize}

\role{Undergraduate Researcher, VCL Lab, Seoul National University of Science and Technology}{Two-stage parcel inventory recognition}{Jan 2023 - Dec 2023}
\begin{itemize}
  \item Built two-stage recognition for stacked parcels: DINO and YOLO detect product regions,
        Weighted Boxes Fusion merges them into cleaner boxes than either alone, then EfficientNet-B7
        with test-time augmentation classifies each crop.
  \item Placed 1st of 19 teams, identifying 56.6\% of the 53 evaluation products against 49.1\% for
        the runner-up, at 0.80 mean IoU.
\end{itemize}

%==============================================================
% PROJECTS
%==============================================================
\section{Selected Projects}
\textbf{Road Damage Analysis} --- Pipeline that pulls road geometry from OpenStreetMap and the
matching Street View frames, detects pavement damage with YOLO, and renders an interactive map with
clustering and a heatmap for triage.\\[1pt]
\textbf{Marionette Motion Control} --- Kinect RGB-D 3D joint capture driving a C++ servo control layer
on an embedded controller, each servo clamped to its reachable range so a bad pose estimate cannot
damage the rig.

%==============================================================
% PUBLICATIONS, AWARDS, AND PRESENTATIONS
%==============================================================
\section{Publications, Awards, and Presentations}
\textit{Vision-Based Construction Worker Safety Monitoring with Pose-Derived Workload and
Ergonomic-Risk Proxies}\\
Under review, Automation in Construction \hfill 2026\\[2pt]
\textbf{Research Grant}, BuiltWorlds Global Summit (\$12,500) \hfill 2026\\
\textbf{1st Prize}, CJ Logistics Future Technology Challenge (\$15,000) \hfill 2023\\[2pt]
\textit{Sensor-Free Wearable System for Activity and Fatigue Detection in Worker Safety Monitoring},
BuiltWorlds Global Summit \hfill 2026\\
\textit{Real-Time Control of Microfluidic Mixing via Machine Learning-Driven Prediction},
BMES Annual Meeting \hfill 2025\\
\textit{Development of Marionette Motion Control System}, ICROS-KROS Conference \hfill 2022

\end{document}
